\documentclass[10pt]{article}

\usepackage[utf8]{inputenc}

\usepackage[T1]{fontenc}

\usepackage{amsmath}

\usepackage{amsfonts}

\usepackage{amssymb}

\usepackage[version=4]{mhchem}

\usepackage{stmaryrd}

\usepackage{bbold}

\usepackage{graphicx}

\usepackage[export]{adjustbox}

\graphicspath{ {./images/} }



\begin{document}

справедливо п р е д с т а в л е н и е Р и с с а функции $u(x)$ в виде суммы потенциала и гармонич. функции $h(x)$ :



\[

u(x)=-\int_{K} E_{n}(|x-y|) d \mu(y)+h(x),

\]



где



\[

E_{2}(|x-y|)=\ln \frac{1}{|x-y|}, E_{n}(|x-y|)=\frac{1}{|x-y|^{n-2}},

\]



$n \geqslant 3,|x-y|$ - расстояние между точками $x, y \in \mathbb{R}^{n}$ (см. [1]). Мера $\mu$ наз. а с с о ц и и р о в а н н ой м е$\mathrm{p}$ о й для функции $u(x)$ или м е $\mathrm{p}$ о й $\mathrm{P}$ и с с $\mathrm{a}$.



Если $K=\bar{H}$ есть замыкание области $H$, причем существует обобщенная функция Грина $g(x, y ; H)$, то формулу (1) можно записать в виде



\[

u(x)=-\int_{\bar{H}} g(x, y ; H) d \mu(y)+h^{*}(x),

\]



где $h^{*}(x)$ - наименьшая гармонич. мажоранта $u(x)$ в области $H$.



Формулы (1), (2) можно распространить при нек-рых дополнительных условиях на всю область $D$ (см. $C y 6-$ гармоническая функция, а также [3], [5]).



\begin{enumerate}

  \setcounter{enumi}{1}

  \item Р.т. о с реднем значени и с уб гармон и ч е с к о й ф у нк ц и и: если $u(x)$-субгармонич. функция в кольцевой области $\left\{x \in \mathbb{R}^{n}: 0 \leqslant r \leqslant \mid x-\right.$ $\left.-x_{0} \mid \leqslant R\right\}$, то ее среднее значение по площади сферы $S_{n}\left(x_{0}, \rho\right)$ с центром $x_{0}$ и радиусом $\rho, r \leqslant \rho \leqslant R$, равно

\end{enumerate}



\[

J(\rho)=J\left(\rho ; x_{0}, u\right)=\frac{1}{\sigma_{n}(\rho)} \int_{S_{n}\left(x_{0}, \rho\right)} u(y) d \sigma_{n}(y),

\]



где $\sigma_{n}(\rho)$ - площадь $S_{n}\left(x_{0}, \rho\right)$, и является выпуклой функцией относительно $1 / \rho^{n-2}$ при $n \geqslant 3$ и относительно $\ln \rho$ при $n=2$. Если же $u(x)-$ субгармонич. функция во всем шаре $\left\{x \in \mathbb{R}^{n}:\left|x-x_{0}\right| \leqslant R\right\}$, то $J(\rho)$, кроме того, - неубывающая непрерывная функция относительно $\rho$ при условии, что $J(0)=u\left(x_{0}\right)$ (cм. [1]).



\begin{enumerate}

  \setcounter{enumi}{2}

  \item Р.т. о б аналитических функция кл а с с о в Х а р д и $H_{\delta}, \delta>0$ : если $f(z)-$ регулярная аналитич. функция в единичном круге $D=$ $=\left\{z=r e^{i \theta} \in \mathbb{C}:|z|<1\right\}$ класса Харди $H_{\delta}, \delta>0$ (см. Граничные свойства аналитических функций), то для нее имеют место соотношения

\end{enumerate}



\[

\begin{gathered}

\lim _{r \rightarrow 1} \int_{E}\left|f\left(r e^{i \theta}\right)\right|^{\delta} d \theta=\int_{E}\left|f\left(e^{i \theta}\right)\right|^{\delta} d \theta, \\

\lim _{r \rightarrow 1} \int_{0}^{2 \pi}\left|f\left(r e^{i \theta}\right)-f\left(e^{i \theta}\right)\right|^{\delta} d \theta=0,

\end{gathered}

\]



где $E$ - любое множество положительной меры на окружности $\Gamma=\left\{z=e^{i \theta}:|z|=1\right\}, \quad f\left(e^{i \theta}\right)-$ граничные значения $f(z)$ на $Г$. Кроме того, $f(z) \in H_{1}$ тогда и только тогда, когда ее первообразная непрерывна в замкнутом круге $D \cup \Gamma$ и абсолютно непрерывна на $\Gamma$ (см. [2]).



Теоремы 1) - 3) доказаны Ф. Риссом (см. [1], [2]). Jum.: [1] R i e s z F.. "Acta math.", 1926, v. 48, p. 329-43; 1930 , v. 54, p. 321-60; [2] e r o ж e, "Math. Z.», 1923, Bd 18, S. $87-95 ;$ [3] П р и в а л о в И. И., Субгармонические финукции, М. -Л., 1937; [4] е го ж е, Граничные свойства аналитических функций, 2 изд., М.-Ј., 1950; [5] X е й м а н У. К., К е н н е д и П. Б., Субгармонические функции, пер. с англ., т. 1, м., 1980 . Б., Субгармонические функции, пер. с англ.,



\includegraphics[max width=\textwidth, center]{2023_07_08_52dff718f3c203cc5f97g-1(2)}

$\ln M(\alpha, \beta)$ точной верхней грани модуля $M(\alpha, \beta)$ билинейной формы



\[

\sum_{i=1}^{m} \sum_{j=1}^{n} a_{i j} x_{i} y_{j}

\]



на множестве



\[

\sum_{i=1}^{m}\left|x_{i}\right|^{1 / \alpha} \leqslant 1, \sum_{j=1}^{m}\left|y_{j}\right|^{1 / \beta} \leqslant 1

\]



(если $\alpha=0$ или $\beta=0$, то соответственно $\left|x_{i}\right| \leqslant 1, i=1, \ldots$ .., m или $\left.\left|y_{j}\right| \leqslant 1, j=1, \ldots, n\right)$ является выпуклой Функцией от параметров $\alpha$ и $\beta$ в области $\alpha \geqslant 0, \beta \geqslant 0$, если форма действительна $\left(a_{i j}, x_{i}, y_{i} \in \mathbb{R}+\right)$, и в области $0 \leqslant \alpha, \quad \beta \leqslant 1, \quad \alpha+\beta \geqslant 1$, если форма комплексна $\left(a_{i j}, x_{i}, y_{j} \in \mathbb{C}\right)$. Эта теорема была доказана М. РисОбобщение Р. т. в. на линейные операторы (см. [3]): пусть $L_{p}, 1 \leqslant p \leqslant \infty,-$ совокупность всех комплексозначных суммируемых в $p$-й степени при $1 \leqslant p<\infty$ и существенно ограниченных при $p=\infty$ функций на нек-ром пространстве с мерой; пусть, далее, $T: L_{p_{i}} \rightarrow$ $\rightarrow L_{q_{i}}, 1 \leqslant p_{i}, \quad q_{i} \leqslant \infty, i=0,1,-$ линейный непрерывный оператор; тогда $T$ является непрерывным оператором из $L_{p_{t}}$ в $L_{q_{t}}$, где



$1 / p_{t}=(1-t) / p_{0}+t / p_{1}, 1 / q_{t}=(1-t) / q_{0}+t / q_{1}, t \in[0,1]$



и норма $k_{t}$ оператора $T$ (как оператора из $L_{p_{t}}$ в $L_{q_{t}}$ ) удовлетворяет неравенству $k_{t} \leqslant k_{0}^{1-t} k_{1}^{t}$ (т. е. является логарифмически выпуклой функцией). Әту теорему наз. т е о р е м о й $\mathbf{P}$ и с с $\mathbf{a}-\mathrm{T} о р и$ и о о б инт е р иоляции, но иногда те о ре мой выII укл о с т и $\mathbf{Р}$ ис с а [4].



Р. т. в. явилась отправным пунктом для целого направления в анализе, где изучаются интерџоляционные свойства линейных операторов. Среди первых из обобщений Р. т. в.- т е о р е м а $\mathrm{M}$ а р ц и н к е в и ч а об интерполяции [5], к-рая гарантирует при $1 \leqslant p_{i} \leqslant$ $\leqslant q_{i} \leqslant \infty, i=0,1$, непрерывность оператора $T: L_{p_{t}} \rightarrow$ $\rightarrow L_{q_{t}}, \quad t \in(0,1)$ при более слабых предположениях, чем в теореме Рисса - Торина. См. также Интерполирование операторов.



Jum.: [1] R i e s Z M., "Acta math.", 1926, v. 49, p. 465-97; [2] Х а р д и Г., Ли т т Лу в в у Д Д., П о ли а Г., Неравенcтва, пер. c aнrл., M., 1948; [3] T h o r i n G., "Comm. Sém. $\Gamma$., Введение в гармонический анализ на евклидовых пространствах, пер. с англ., M., 1974; [5] M a r c i n k i e w i c z J., "C. r. Acad. Sci.», 1939, t. 208, p. 1272-73; [6] К p е й н C. Г., нейных операторов, м., 1978; [7] T $\mathrm{T}$ i e b e 1 H., Interpolation theory, B., 1978 . $\quad$ B. М. Тихомиров.



РИССА - ФИІІЕРА ТЕОРЕМА - теорема, устанавливающая связь между пространствами $l^{2}$ и $L^{2}[a, b]$ : если система функциї $\left\{\varphi_{n}(t)\right\}_{n=1}^{\infty}$ ортонормирована на отрезке $[a, b]$, а последовательность чисел $\left\{c_{n}\right\}_{n=1}^{\infty}$ такова, что



\[

\sum_{n=1}^{\infty} c_{n}^{2}<\infty

\]



(то есть $c_{n} \in l^{2}$ ), то существует функция $f(t) \in L^{2}[a, b]$, для к-рой



\[

\int_{a}^{b}|f(t)|^{2} d t=\sum_{n=1}^{\infty} c_{n}^{2}, c_{n}=\int_{a}^{b} f(t) \varphi_{n}(t) d t .

\]



При этом функция $f(t)$ единственна как әлемент пространства $L^{2}[a, b]$, т. е. с точностью до ее значений на множестве нулевой меры Лебега. В частности, если ортонормированная система $\left\{\varphi_{n}(t)\right\}$ замкнута (полна) в $L^{2}[a, b]$, то с помощью P. - Ф. т. устанав.іивается, что пространства $l^{2}$ и $L^{2}[a, b]$ изоморфны и изометричны.



Р. - Ф. т. доказана независимо Ф. Риссом [1] и Э. Фишером [2].



Jum.: [1] R i e s z F., "C. r. Acad. sci.", 1907, t. 144, p. 615-



\begin{center}

\includegraphics[max width=\textwidth]{2023_07_08_52dff718f3c203cc5f97g-1(1)}

\end{center}



\begin{center}

\includegraphics[max width=\textwidth]{2023_07_08_52dff718f3c203cc5f97g-1}

\end{center}



РИССОВ ТЕОРЕМА - 1) Р. т. е д и н с т в е н нос ти для ограниченных аналитическ и х ф у к ц и й: если $f(z)$ - ограниченная регулярная аналитич. функция в единичном круге $D=\{z \in \mathbb{C}$ : $|z|<1\}$, имеющая радиальные граничные значения нуль на множестве $E$ точек окружности $\Gamma=\{z:|z|=1\}$ положительной меры, mes $E>0$, то $f(z) \equiv 0$. Теорема





\end{document}